\chapter{Capitolo 14 - Governance epistemica: versioning, audit trail, quality gates}

\section{1. Problema}

Una teoria può avere buone definizioni e buoni benchmark, ma fallire comunque come programma scientifico se non possiede una governance epistemica esplicita.

Nel contesto OCT, il rischio non e solo tecnico. E istituzionale:

\begin{enumerate}
\def\labelenumi{\arabic{enumi}.}
\tightlist
\item
  stati dei claim aggiornati in modo incoerente tra file diversi;
\item
  release pubbliche prive di tracciabilità decisionale;
\item
  conflitti tra risultati locali e narrativa generale del progetto.
\end{enumerate}

Per questo il problema del capitolo e netto:

\textbf{quale sistema di regole permette di governare nel tempo il passaggio dei claim tra \texttt{draft}, \texttt{in\_proof}, \texttt{validated/revise/reject}, mantenendo trasparenza, reversibilità controllata e responsabilità editoriale?}

\subsection{1.1 Perché la governance e parte della teoria}

In OCT la validità non e solo proprietà logica di una formula: e proprietà processuale di una catena evidenziale.

Di conseguenza:

\begin{enumerate}
\def\labelenumi{\arabic{enumi}.}
\tightlist
\item
  senza governance, la teoria non scala oltre il laboratorio originario;
\item
  senza audit trail, la validazione non e verificabile retrospettivamente;
\item
  senza quality gate, la pubblicazione rischia di confondere stato ipotetico e stato consolidato.
\end{enumerate}

\subsection{\texorpdfstring{1.2 Ruolo nella milestone \texttt{v1.4}}{1.2 Ruolo nella milestone v1.4}}

Con questo capitolo si chiude il blocco metodologico.

\begin{enumerate}
\def\labelenumi{\arabic{enumi}.}
\tightlist
\item
  Capitolo 12: protocollo di decisione;
\item
  Capitolo 13: infrastruttura benchmark;
\item
  Capitolo 14: costituzione operativa del corpus.
\end{enumerate}

Dopo questo punto il libro può passare ai case study (Capitoli 15-16) senza perdere rigore di stato.

\begin{center}\rule{0.5\linewidth}{0.5pt}\end{center}

\section{2. Tesi del capitolo}

La tesi e articolata in dieci enunciati.

\begin{enumerate}
\def\labelenumi{\arabic{enumi}.}
\tightlist
\item
  Ogni claim OCT deve avere doppio stato: formale e decisionale.
\item
  Lo stato decisionale e valido solo se ancorato a evidenza tracciata e versionata.
\item
  I quality gate devono essere pubblici, finiti e verificabili ex post.
\item
  Ogni cambio di stato deve essere reversibile solo tramite procedura esplicita.
\item
  Le deviazioni dal protocollo sono ammesse solo con eccezione registrata.
\item
  L'audit trail e parte del risultato scientifico, non metadato accessorio.
\item
  Il versioning epistemico deve distinguere evoluzione teorica da evoluzione empirica.
\item
  La pubblicazione multi-piattaforma richiede sincronizzazione di stato e identificatori.
\item
  Le decisioni \texttt{revise} non sono fallimenti: sono meccanismi di controllo della deriva.
\item
  Il sistema proposto e sufficiente per governare la fase case-study e la transizione a preprint stabile.
\end{enumerate}

\begin{center}\rule{0.5\linewidth}{0.5pt}\end{center}

\section{3. Definizioni canoniche}

\subsection{Definizione 14.1 (Stato formale)}

Stato formale di un claim nel programma teorematico:

\texttt{S\_form\ in\ \{draft,\ in\_proof,\ validated\}}.

\subsection{Definizione 14.2 (Stato decisionale)}

Stato decisionale basato su evidenza operativa:

\texttt{S\_dec\ in\ \{validated,\ revise,\ reject\}}.

\texttt{S\_dec} non sostituisce \texttt{S\_form}, ma lo qualifica per uso editoriale e pubblicazione.

\subsection{Definizione 14.3 (Record di stato)}

Un record minimo di stato e:

\texttt{R\_state\ :=\ \textless{}ClaimID,\ S\_form,\ S\_dec,\ EvidenceRef,\ DecisionDate,\ Version,\ Signatures\textgreater{}}.

\subsection{Definizione 14.4 (Quality Gate)}

Un quality gate \texttt{G\_k} e una condizione booleana preregistrata che deve risultare \texttt{PASS} per autorizzare un tipo specifico di rilascio.

\subsection{Definizione 14.5 (Eccezione governance)}

Eccezione governance: deviazione motivata da procedura standard, registrata con:

\begin{enumerate}
\def\labelenumi{\arabic{enumi}.}
\tightlist
\item
  motivo;
\item
  rischio introdotto;
\item
  mitigazione;
\item
  data di scadenza dell'eccezione.
\end{enumerate}

\subsection{Definizione 14.6 (Audit trail epistemico)}

Catena ordinata di evidenze e decisioni che permette di ricostruire integralmente:

\begin{enumerate}
\def\labelenumi{\arabic{enumi}.}
\tightlist
\item
  quali dati sono stati usati;
\item
  quali criteri erano attivi;
\item
  perché uno stato e cambiato.
\end{enumerate}

\begin{center}\rule{0.5\linewidth}{0.5pt}\end{center}

\section{4. Sviluppo formale}

\subsection{4.1 Architettura governance a quattro strati}

La governance proposta e composta da quattro strati coordinati.

\begin{enumerate}
\def\labelenumi{\arabic{enumi}.}
\tightlist
\item
  \textbf{Strato G1 - Semantico}: definizioni, notazione, perimetro del claim.
\item
  \textbf{Strato G2 - Evidenziale}: benchmark, metriche, report, repliche.
\item
  \textbf{Strato G3 - Decisionale}: matrice \texttt{validated/revise/reject}.
\item
  \textbf{Strato G4 - Rilascio}: gate, versioning, pubblicazione.
\end{enumerate}

Un errore in un singolo strato può invalidare il rilascio anche se gli altri tre sono coerenti.

\subsection{4.2 Doppio stato e regole di coerenza}

Definiamo regole minime di coerenza tra \texttt{S\_form} e \texttt{S\_dec}.

\begin{enumerate}
\def\labelenumi{\arabic{enumi}.}
\tightlist
\item
  se \texttt{S\_form\ =\ draft}, allora \texttt{S\_dec} non può essere \texttt{validated};
\item
  se \texttt{S\_dec\ =\ reject}, il claim non può essere promosso senza nuova evidenza versionata;
\item
  se \texttt{S\_dec\ =\ revise}, il claim resta pubblicabile solo con etichetta esplicita di provvisorieta;
\item
  promozione a \texttt{S\_form\ =\ validated} richiede allineamento con decisione non conflittuale e gate di rilascio superati.
\end{enumerate}

Queste regole evitano inversioni narrative del tipo ``teorema consolidato'' con evidenza ancora in revisione.

\subsection{4.3 Versioning epistemico}

Il versioning deve distinguere almeno tre dimensioni:

\begin{enumerate}
\def\labelenumi{\arabic{enumi}.}
\tightlist
\item
  \texttt{V\_theory}: cambi nella formalizzazione/assiomatica;
\item
  \texttt{V\_evidence}: cambi nei benchmark, dataset o analisi;
\item
  \texttt{V\_release}: confezionamento editoriale/pubblicazione.
\end{enumerate}

Regola:

\begin{enumerate}
\def\labelenumi{\arabic{enumi}.}
\tightlist
\item
  una modifica in \texttt{V\_evidence} senza modifica teorica non autorizza riscrittura implicita di claim;
\item
  una modifica in \texttt{V\_theory} richiede rivalutazione esplicita delle evidenze collegate;
\item
  ogni release pubblica deve dichiarare la tripla versione (\texttt{V\_theory}, \texttt{V\_evidence}, \texttt{V\_release}).
\end{enumerate}

\subsection{\texorpdfstring{4.4 Quality gate standard \texttt{G0-G6}}{4.4 Quality gate standard G0-G6}}

Proponiamo sette gate minimi.

\begin{enumerate}
\def\labelenumi{\arabic{enumi}.}
\tightlist
\item
  \texttt{G0} - completezza metadati (\texttt{PASS/FAIL});
\item
  \texttt{G1} - integrità artefatti (manifest + hash);
\item
  \texttt{G2} - preregistrazione e criteri decisionali bloccati;
\item
  \texttt{G3} - evidenza benchmark indipendente minima;
\item
  \texttt{G4} - audit trail leggibile e ripercorribile;
\item
  \texttt{G5} - coerenza tra stato formale e decisionale;
\item
  \texttt{G6} - checklist pubblicazione multi-piattaforma allineata.
\end{enumerate}

Policy:

\begin{enumerate}
\def\labelenumi{\arabic{enumi}.}
\tightlist
\item
  \texttt{FAIL} su \texttt{G0-G2} blocca qualsiasi release;
\item
  \texttt{FAIL} su \texttt{G3-G5} blocca release ``claim-strong'' e consente solo release di lavoro;
\item
  \texttt{FAIL} su \texttt{G6} consente rilascio locale ma non rilascio ufficiale.
\end{enumerate}

\subsection{4.5 Procedura di cambio stato (PCS-6)}

Ogni cambio stato deve seguire sei passi.

\begin{enumerate}
\def\labelenumi{\arabic{enumi}.}
\tightlist
\item
  proposta di cambio con motivazione;
\item
  verifica automatica gate disponibili;
\item
  revisione umana indipendente minima;
\item
  emissione decisione e firma;
\item
  aggiornamento record \texttt{R\_state};
\item
  pubblicazione del changelog decisionale.
\end{enumerate}

La PCS-6 rende il cambio di stato una procedura, non un atto implicito.

\subsection{4.6 Gestione eccezioni e debito epistemico}

Le eccezioni non vanno eliminate, vanno governate.

Per ogni eccezione registriamo:

\begin{enumerate}
\def\labelenumi{\arabic{enumi}.}
\tightlist
\item
  grado di rischio (\texttt{low/medium/high});
\item
  scadenza entro cui rientrare in procedura standard;
\item
  owner responsabile del rientro;
\item
  impatto su gate e release ammesse.
\end{enumerate}

L'accumulo di eccezioni oltre soglia definisce debito epistemico e obbliga una release dedicata di consolidamento.

\subsection{4.7 Audit trail minimo pubblicabile}

Ogni rilascio claim-sensitive deve includere:

\begin{enumerate}
\def\labelenumi{\arabic{enumi}.}
\tightlist
\item
  decision matrix corrente;
\item
  stato teoremi allineato;
\item
  riferimenti ai report benchmark rilevanti;
\item
  log sintetico dei cambi di stato;
\item
  elenco eccezioni aperte e chiuse.
\end{enumerate}

Senza questo bundle, un revisore esterno non può verificare la storia decisionale.

\subsection{4.8 Governance multi-piattaforma}

Per GitHub, Zenodo, OSF, Hugging Face il principio e uno: \textbf{stesso stato, stessa versione, stesso claim ledger}.

Sincronizzazioni obbligatorie:

\begin{enumerate}
\def\labelenumi{\arabic{enumi}.}
\tightlist
\item
  ID versione release;
\item
  checksum dei pacchetti dati;
\item
  decision table di riferimento;
\item
  timestamp di emissione.
\end{enumerate}

Divergenze tra piattaforme sono trattate come incidenti di governance.

\subsection{4.9 Politica di comunicazione dello stato}

Ogni documento pubblico che cita un claim deve riportare etichetta di stato nel punto d'uso.

Formato raccomandato:

\texttt{ClaimID\ {[}S\_form\ /\ S\_dec\ /\ EvidenceLevel{]}}.

Esempio:

\texttt{D03\ {[}in\_proof\ /\ revise\ /\ S2{]}}.

Questo riduce ambiguità per lettori nuovi e impedisce overclaiming involontario.

\subsection{4.10 Condizione di sufficienza del capitolo}

La governance e sufficiente se garantisce contemporaneamente:

\begin{enumerate}
\def\labelenumi{\arabic{enumi}.}
\tightlist
\item
  tracciabilità completa dei cambi stato;
\item
  regole di blocco rilascio verificabili;
\item
  allineamento stabile tra teoria, evidenza e pubblicazione.
\end{enumerate}

La condizione e soddisfatta dall'integrazione di G1-G4 con PCS-6, gate \texttt{G0-G6} e policy multi-piattaforma.

\subsection{4.11 Indicatori di salute della governance (KPI-E)}

Per monitorare se la governance funziona nel tempo proponiamo cinque indicatori periodici:

\begin{enumerate}
\def\labelenumi{\arabic{enumi}.}
\tightlist
\item
  \texttt{K1\_time\_to\_decision}: tempo medio tra chiusura benchmark e decisione ufficiale;
\item
  \texttt{K2\_trace\_completeness}: percentuale record con audit trail completo;
\item
  \texttt{K3\_exception\_debt}: numero eccezioni aperte oltre scadenza;
\item
  \texttt{K4\_sync\_gap}: discrepanze di stato tra piattaforme pubbliche;
\item
  \texttt{K5\_revise\_stagnation}: claim in \texttt{revise} senza piano attivo di uscita.
\end{enumerate}

Uso operativo:

\begin{enumerate}
\def\labelenumi{\arabic{enumi}.}
\tightlist
\item
  soglie KPI fissate trimestralmente;
\item
  superamento soglia produce azione correttiva obbligatoria;
\item
  trend peggiorativo su due finestre consecutive blocca release non critiche.
\end{enumerate}

Questo blocco evita che la governance resti un disegno statico e la trasforma in sistema vivo, misurabile e migliorabile.

\subsection{4.12 Registro delle decisioni contestate}

Una governance matura deve prevedere anche disaccordo strutturato.\\
Per questo ogni contestazione formale a una decisione va registrata in un log dedicato con:

\begin{enumerate}
\def\labelenumi{\arabic{enumi}.}
\tightlist
\item
  decisione contestata e riferimento al record \texttt{R\_state};
\item
  motivo tecnico della contestazione;
\item
  evidenza addizionale proposta;
\item
  esito della revisione (conferma, modifica, sospensione).
\end{enumerate}

Il registro non indebolisce l'autorità del processo: la rende verificabile e correggibile.

\begin{center}\rule{0.5\linewidth}{0.5pt}\end{center}

\section{5. Proposizioni e teoremi locali}

\subsection{Proposizione 14.1 (Necessità del doppio stato)}

Enunciato:
Un singolo stato non distingue adeguatamente maturità formale e maturità evidenziale del claim.

Intuizione:
la stessa formula può essere formalmente elegante ma sperimentalmente ancora instabile.

Stato: \texttt{validated}.

\subsection{Proposizione 14.2 (Funzione anti-deriva dei gate)}

Enunciato:
Un sistema di gate pubblici riduce la probabilità di promozioni opportunistiche di stato.

Intuizione:
vincoli espliciti rendono costoso ignorare evidenze negative o incomplete.

Stato: \texttt{validated}.

\subsection{Proposizione 14.3 (Valore epistemico delle eccezioni tracciate)}

Enunciato:
Eccezioni registrate e temporalmente limitate migliorano l'affidabilità complessiva rispetto a eccezioni implicite non documentate.

Intuizione:
la trasparenza dell'anomalia e preferibile alla sua invisibilità.

Stato: \texttt{in\_review}.

\subsection{Teorema 14.4 (Consistenza procedurale della PCS-6)}

Ipotesi:

\begin{enumerate}
\def\labelenumi{\arabic{enumi}.}
\tightlist
\item
  applicazione completa della procedura PCS-6;
\item
  quality gate valutati prima della decisione;
\item
  aggiornamento record di stato obbligatorio.
\end{enumerate}

Tesi:
La governance produce cambi di stato auditabili e confrontabili nel tempo.

Proof sketch:

\begin{enumerate}
\def\labelenumi{\arabic{enumi}.}
\tightlist
\item
  la procedura serializza decisione e responsabilità;
\item
  i gate impongono condizioni minime verificabili;
\item
  il record standardizzato abilità confronto storico.
\end{enumerate}

Conseguenze:
rende possibile una storia epistemica computabile del corpus OCT.

Stato: \texttt{in\_proof}.

\subsection{\texorpdfstring{Teorema 14.5 (Sufficienza dei gate \texttt{G0-G6} per rilascio claim-sensitive)}{Teorema 14.5 (Sufficienza dei gate G0-G6 per rilascio claim-sensitive)}}

Ipotesi:

\begin{enumerate}
\def\labelenumi{\arabic{enumi}.}
\tightlist
\item
  tutti i gate \texttt{G0-G6} in \texttt{PASS};
\item
  allineamento tra ledger, decision matrix e stato teoremi;
\item
  assenza di eccezioni \texttt{high} scadute.
\end{enumerate}

Tesi:
Il rilascio e metodologicamente idoneo a disseminazione scientifica pubblica.

Proof sketch:

\begin{enumerate}
\def\labelenumi{\arabic{enumi}.}
\tightlist
\item
  \texttt{G0-G2} garantiscono base tecnica minima;
\item
  \texttt{G3-G5} garantiscono coerenza evidenziale e semantica;
\item
  \texttt{G6} garantisce consistenza cross-piattaforma.
\end{enumerate}

Conseguenze:
fornisce criterio operativo unificato per pubblicazione OCT.

Stato: \texttt{in\_proof}.

\begin{center}\rule{0.5\linewidth}{0.5pt}\end{center}

\section{6. Implicazioni operative}

\subsection{\texorpdfstring{6.1 Per chiusura milestone \texttt{v1.4}}{6.1 Per chiusura milestone v1.4}}

Con il Capitolo 14 la milestone metodologica e chiusa quando:

\begin{enumerate}
\def\labelenumi{\arabic{enumi}.}
\tightlist
\item
  Capitolo 12 (protocollo) e completo;
\item
  Capitolo 13 (benchmark design) e completo;
\item
  Capitolo 14 (governance) e completo e coerente con gli artefatti di validazione.
\end{enumerate}

\subsection{6.2 Per i case study (Capitoli 15-16)}

Ogni case study dovra essere scritto con vincolo di stato esplicito:

\begin{enumerate}
\def\labelenumi{\arabic{enumi}.}
\tightlist
\item
  quali claim usa;
\item
  in quale stato sono;
\item
  quali parti sono ipotesi operative e quali sono evidenze consolidate.
\end{enumerate}

\subsection{6.3 Per gestione editoriale}

La governance permette un editing ``forte'' senza alterare significato epistemico:

\begin{enumerate}
\def\labelenumi{\arabic{enumi}.}
\tightlist
\item
  si può migliorare il testo;
\item
  non si può migliorare lo stato senza nuova evidenza.
\end{enumerate}

\subsection{6.4 Per revisori e comunità scientifica}

Un revisore esterno può ricostruire l'intera catena decisionale senza accesso privato al team, se il bundle governance e completo.

Questa proprietà e essenziale per credibilità internazionale del programma OCT.

\subsection{6.5 Per ritmo operativo del team}

La governance non deve solo "controllare"; deve anche dare ritmo alla ricerca.

Cadenza consigliata:

\begin{enumerate}
\def\labelenumi{\arabic{enumi}.}
\tightlist
\item
  review settimanale dei gate per i claim attivi;
\item
  review mensile del debito epistemico;
\item
  review trimestrale dei KPI-E e delle policy eccezioni;
\item
  retrospettiva semestrale su versioning e coerenza multi-piattaforma.
\end{enumerate}

Questa cadenza rende sostenibile il programma su orizzonte pluriennale senza accumulare disordine metodologico.

\begin{center}\rule{0.5\linewidth}{0.5pt}\end{center}

\section{7. Limiti e non-claim}

Limiti espliciti:

\begin{enumerate}
\def\labelenumi{\arabic{enumi}.}
\tightlist
\item
  il modello di governance e baseline operativa, non costituzione definitiva;
\item
  alcuni domini ad alta sensibilità possono richiedere gate aggiuntivi;
\item
  la procedura non elimina il conflitto interpretativo, ma lo rende tracciabile.
\end{enumerate}

Failure mode riconosciuti:

\begin{enumerate}
\def\labelenumi{\arabic{enumi}.}
\tightlist
\item
  proliferazione di eccezioni non chiuse;
\item
  aggiornamento parziale dei record tra piattaforme;
\item
  uso improprio dello stato \texttt{revise} come zona grigia permanente;
\item
  eccesso di burocrazia che rallenta la ricerca viva.
\end{enumerate}

Contromisure raccomandate:

\begin{enumerate}
\def\labelenumi{\arabic{enumi}.}
\tightlist
\item
  revisione periodica del debito epistemico;
\item
  allineamento automatico dei ledger principali;
\item
  policy di scadenza massima per claim in \texttt{revise} senza piano;
\item
  semplificazione continua dei template senza perdere auditabilità.
\end{enumerate}

Box C - Non-claim

Questo capitolo non implica:

\begin{enumerate}
\def\labelenumi{\arabic{enumi}.}
\tightlist
\item
  che la governance sostituisca il valore creativo della teoria;
\item
  che ogni decisione governance sia incontestabile;
\item
  che un claim in \texttt{revise} sia scientificamente irrilevante.
\end{enumerate}

\begin{center}\rule{0.5\linewidth}{0.5pt}\end{center}

\section{8. Claim Ledger (S0-S3)}

{\def\LTcaptype{none} % do not increment counter
\begin{longtable}[]{@{}lllll@{}}
\toprule\noalign{}
ID & Claim & Tipo & Evidenza & Stato \\
\midrule\noalign{}
\endhead
\bottomrule\noalign{}
\endlastfoot
C14.1 & Il doppio stato (\texttt{S\_form}, \texttt{S\_dec}) e necessario per evitare ambiguità epistemiche & governance-metateorica & S1 & validated \\
C14.2 & I gate \texttt{G0-G6} riducono la deriva decisionale in rilascio & metodologico-operativo & S1 & validated \\
C14.3 & La PCS-6 rende auditabili i cambi di stato & procedurale & S2 & in\_proof \\
C14.4 & La sincronizzazione multi-piattaforma e condizione di credibilità esterna & infrastrutturale & S1 & in\_review \\
C14.5 & La gestione esplicita delle eccezioni riduce il rischio epistemico cumulativo & governance-rischio & S1 & in\_review \\
C14.6 & Il pacchetto governance e sufficiente per supportare la fase case-study & metaprogramma & S2 & in\_proof \\
\end{longtable}
}

\begin{center}\rule{0.5\linewidth}{0.5pt}\end{center}

\section{9. Bridge al Capitolo 15}

Con il Capitolo 14 il quadro metodologico dell'opera e completo:

\begin{enumerate}
\def\labelenumi{\arabic{enumi}.}
\tightlist
\item
  protocollo di validazione (Cap. 12);
\item
  benchmark design e replicabilità (Cap. 13);
\item
  governance di stato e rilascio (Cap. 14).
\end{enumerate}

Il Capitolo 15 entra nel primo caso studio fisico (standing waves e controfase strutturale), usando integralmente questa architettura come base di lettura e valutazione.

In sintesi: da qui in avanti l'opera non espone solo teoria OCT, ma mostra OCT in azione con disciplina scientifica esplicita.

\begin{center}\rule{0.5\linewidth}{0.5pt}\end{center}

\section{Riferimenti}

Riferimenti di framework interno:

\begin{enumerate}
\def\labelenumi{\arabic{enumi}.}
\tightlist
\item
  Ghioni, F. (2026). \texttt{TE\_CORE\_v5.1.md}.
\item
  Ghioni, F. (2026). \texttt{Ordinative\_Set\_Theory\_OST\_A\_Concise\_Guide\_For\_AI\_v2\_1.md}.
\item
  \texttt{OCT\_FOUNDATIONAL\_BOOK\_CHAPTER\_12\_v1\_0.md}.
\item
  \texttt{OCT\_FOUNDATIONAL\_BOOK\_CHAPTER\_13\_v1\_0.md}.
\item
  \texttt{OCT\_Theory\_and\_Theorems/validation/DECISION\_MATRIX\_FINAL\_UNIFIED\_v0\_1.md}.
\item
  \texttt{OCT\_Theory\_and\_Theorems/validation/OCT\_THEOREM\_STATUS\_ALIGNMENT\_v1\_0\_candidate\_v0\_1.md}.
\item
  \texttt{OCT\_Theory\_and\_Theorems/validation/OCT\_PUBLICATION\_PROGRESS\_TRACKER\_v0\_1.md}.
\end{enumerate}

Riferimenti primari:

\begin{enumerate}
\def\labelenumi{\arabic{enumi}.}
\tightlist
\item
  Popper, K. (1959). \emph{The Logic of Scientific Discovery}.
\item
  Lakatos, I. (1978). \emph{The Methodology of Scientific Research Programmes}.
\item
  Merton, R. K. (1973). \emph{The Sociology of Science}.
\item
  Goodman, S., Fanelli, D., Ioannidis, J. (2016). \emph{What does research reproducibility mean?}.
\end{enumerate}
